\documentclass[10pt,a4paper,roman]{moderncv}
\moderncvstyle{banking}
\moderncvcolor{purple}
\nopagenumbers{}
\usepackage[utf8]{inputenc}
\usepackage[T1]{fontenc}
\usepackage{lmodern}
\usepackage[scale=0.95,top=0.65cm,bottom=0.65cm,left=0.8cm,right=0.8cm]{geometry}
\setlength{\hintscolumnwidth}{2.25cm}
\setlength{\footskip}{18pt}
\name{Wahid}{SAMER}
\address{Cergy (95), France}{}{}
\mobile{07 53 10 95 74}
\email{samer.ahmed.wahid@gmail.com}
\social[linkedin]{wahid-samer-86b6b820a}
\social[github]{wahidrt}
\title{Alternance M2 - Quantitative Analyst / Quant Research}
\begin{document}
\makecvtitle
\vspace{-3.2em}
\begin{center}
\parbox{0.97\textwidth}{\centering\small\itshape
Profil orienté \textbf{finance quantitative et modélisation stochastique}, combinant mathématiques avancées, pricing de produits dérivés et implémentation Python. Recherche d'une \textbf{alternance M2 de 12 mois à partir du 21 septembre 2026}. Mobilité : \textbf{France, Luxembourg, Suisse, Belgique}.}
\end{center}
\vspace{0.25em}

\section{Formation}
\cventry{2025 -- 2027}{Master Mathématiques Appliquées à l'Ingénierie Financière}{CY Cergy Paris Université / CY Tech}{Cergy}{}{\textbf{M1 :} probabilités, processus stochastiques 1--2, EDP, simulation stochastique, méthodes numériques avancées, évaluation des actifs contingents.\\
\textbf{M2 2026--2027 :} stochastic calculus, theory of contingent claims, model calibration and simulation, fixed income, séries temporelles et machine learning avec Python.}
\cventry{2021 -- 2025}{Licence Mathématiques et Applications}{Université d'Évry Paris-Saclay}{Évry}{}{Formation solide en mathématiques pures et appliquées : probabilités, mesure et intégration, analyse, algèbre, topologie, analyse numérique et programmation.}

\section{Projets ciblés}
\cventry{2026}{Python / \LaTeX}{Mémoire - Pricing d'options : Black--Scholes et Heston}{Projet M1}{}{Étude du lien EDS--EDP, Itô et Feynman--Kac ; dérivation/pricing de Call et Put, Greeks, méthodes numériques et comparaison avec la volatilité stochastique de Heston.}
\cventry{2026}{Python}{Trading quantitatif adaptatif}{Projet M1 - équipe de 4}{}{Cointégration, VECM, filtre de Kalman, Ornstein--Uhlenbeck et régimes de marché ; conception d'une chaîne de backtesting et d'analyse multi-actifs.}
\cventry{2025}{Excel / VBA}{Pricer d'options vanilles et Greeks}{Projet M1}{}{Black--Scholes, calcul de $\Delta,\Gamma,\Theta,\nu,\rho$ et scénarios de stress sur la volatilité.}

\section{Compétences}
\cvitem{Mathématiques}{Calcul stochastique, probabilités, EDP, Feynman--Kac, Black--Scholes, Heston, méthodes numériques, optimisation.}
\cvitem{Finance}{Produits dérivés, pricing, Greeks, gestion de portefeuille, analyse quantitative.}
\cvitem{Programmation}{Python (NumPy, Pandas, SciPy, Scikit-learn, Matplotlib), C/C++, VBA, SQL ; Git/GitHub, \LaTeX.}
\cvitem{Langues}{Français : courant ; Arabe : langue maternelle ; Anglais : intermédiaire.}

\section{Expérience \& Engagement}
\cventry{Mai -- juin 2025}{Professeur de mathématiques - stagiaire}{Lycée Pierre Mendès France}{Vitrolles}{}{Préparation et transmission de notions de niveau Terminale ; développement de la rigueur, de la pédagogie et de la communication scientifique.}
\cventry{2019}{Technicien aéronautique - stagiaire}{G1 Aviation}{Gap}{}{Assemblage/maintenance dans un environnement exigeant en précision, procédures et sécurité.}
\cvitem{Associatif}{Bénévole aux \textbf{Restaurants du Cœur} pendant près de 3 ans : solidarité, accueil et travail d'équipe.}
\end{document}
